\documentclass[12pt,a4paper]{article}

% ── Packages ─────────────────────────────────────────────────────────────────
\usepackage[margin=1in]{geometry}
\usepackage{hyperref}
\hypersetup{
  colorlinks=true,
  linkcolor=blue!60!black,
  urlcolor=blue!60!black,
  citecolor=blue!60!black,
  pdftitle={IAA Technical Handbook},
  pdfauthor={Project AMP},
}
\usepackage{listings}
\usepackage{xcolor}
\usepackage{booktabs}
\usepackage{longtable}
\usepackage{parskip}
\usepackage{microtype}
\usepackage{titlesec}
\usepackage{fancyhdr}
\usepackage{graphicx}
\usepackage{array}
\usepackage{enumitem}

% ── Listing style ─────────────────────────────────────────────────────────────
\definecolor{codebg}{HTML}{F5F5F5}
\definecolor{codecomment}{HTML}{6A9955}
\definecolor{codestring}{HTML}{CE9178}
\definecolor{codekeyword}{HTML}{569CD6}
\definecolor{codenumber}{HTML}{B5CEA8}

\lstset{
  language=Python,
  backgroundcolor=\color{codebg},
  basicstyle=\ttfamily\footnotesize,
  keywordstyle=\color{codekeyword}\bfseries,
  stringstyle=\color{codestring},
  commentstyle=\color{codecomment}\itshape,
  numberstyle=\tiny\color{gray},
  numbers=left,
  numbersep=8pt,
  breaklines=true,
  breakatwhitespace=true,
  showstringspaces=false,
  frame=single,
  rulecolor=\color{gray!30},
  tabsize=4,
  captionpos=b,
  xleftmargin=15pt,
  xrightmargin=5pt,
  aboveskip=10pt,
  belowskip=10pt,
}

% ── Verbatim style for ASCII diagrams ────────────────────────────────────────
\lstdefinestyle{ascii}{
  language={},
  basicstyle=\ttfamily\small,
  backgroundcolor=\color{codebg},
  frame=single,
  rulecolor=\color{gray!30},
  numbers=none,
  breaklines=true,
}

% ── Header / footer ──────────────────────────────────────────────────────────
\pagestyle{fancy}
\fancyhf{}
\fancyhead[L]{\small Intent Anticipation Agent --- Technical Handbook}
\fancyhead[R]{\small Project AMP}
\fancyfoot[C]{\thepage}
\renewcommand{\headrulewidth}{0.4pt}

% ── Colours ───────────────────────────────────────────────────────────────────
\definecolor{ampblue}{HTML}{1E3A5F}
\definecolor{ampgray}{HTML}{6B7280}

% ── Title ─────────────────────────────────────────────────────────────────────
\title{
  {\color{ampblue}\Huge\bfseries Intent Anticipation Agent}\\[6pt]
  {\color{ampgray}\Large Technical Handbook}\\[10pt]
  {\normalsize Project AMP --- Version 0.2.0}
}
\author{Project AMP Engineering}
\date{April 2026}

% ═════════════════════════════════════════════════════════════════════════════
\begin{document}
\maketitle
\thispagestyle{empty}

\vspace{2em}
\begin{center}
\colorbox{ampblue}{\parbox{0.8\textwidth}{\centering\color{white}\bfseries\large
  AMP $|$ Agentic Microservices Prototyping --- Stage 2 Agent}}
\end{center}

\vspace{2em}
\tableofcontents
\clearpage

% ═════════════════════════════════════════════════════════════════════════════
\subsection{Objective}

The Intent Anticipation Agent (IAA) is a machine-learning inference service that predicts
the single most likely next action for a multi-step enterprise workflow case, returning
a ranked set of candidate actions alongside structured task and timeline data for
downstream consumption.

Within the Project AMP pipeline, IAA occupies \textbf{Stage 2}, running in parallel with
the Risk Radar Agent (RRA) and the Governance Policy Agent (GPA), after the Context
Stitching Agent (CSA) has assembled the case's document, email, and event history.
IAA's output --- an \texttt{ActionRecommendation} object --- is forwarded to the
Delegation Agent (DA), which uses the recommended action, the associated task list, and
the escalation flag to decide who (or which automated channel) should handle the case next.

IAA fills the automation gap between coarse risk scoring (RRA) and routing (DA).
Rule-based routing systems fail when the correct next action depends on the
combination of case history, document completeness, and organisational context.
Humans filling this gap introduce latency of hours to days. IAA closes that gap with
sub-second decisions on the standard path and zero-millisecond decisions on the fast path,
while still routing genuinely ambiguous cases to human review via an explicit escalation flag.

% ═════════════════════════════════════════════════════════════════════════════
\subsection{Problem Statement}

Enterprise workflows --- vendor onboarding, insurance claims, fraud investigations,
contract reviews --- share a structural challenge: at any given moment in a case,
the correct next action is not simply the next step in a fixed flowchart.
It depends on what has already happened (prior event sequence), what is missing
(document gaps, incomplete risk data), and what external policy constraints are active
(risk severity, governance rules). This combinatorial space grows exponentially with
the number of workflow variants, making hard-coded rule engines brittle and expensive to maintain.

Pure rule-based systems fail for two reasons. First, they cannot generalise to new
workflow variants without explicit rule authoring. Second, the conditionals interact:
a missing document may or may not require escalation depending on the case's risk level,
which no static rule set can encode without exponential branching.

Large language models (LLMs) alone are also insufficient. Without access to
historical case patterns, an LLM can only reason from its training data, which does not
include the organisation's specific workflow structure, escalation thresholds, or
actor assignments. LLMs also have no notion of which actions have already been tried,
causing them to repeat or skip steps.

IAA's hybrid approach resolves both limitations:
\begin{enumerate}[nosep]
  \item \textbf{Fast path} --- deterministic rules bypass the LLM entirely for
    unambiguous cases (high risk, missing inputs), achieving \textasciitilde{}0\,ms latency.
  \item \textbf{RAG augmentation} --- the Qdrant vector store retrieves historically
    similar cases (real BPIC event logs and synthetic cases), grounding the LLM prompt
    with observed next-action patterns from analogous workflows.
  \item \textbf{Single-shot LLM chain} --- a merged prompt elicits predicted actions,
    task list, and workflow timeline in one call, cutting the original 3-call design
    from \textasciitilde{}26\,s to \textasciitilde{}8--10\,s.
  \item \textbf{Post-LLM rule overrides} --- hard rules applied after the LLM call
    ensure policy compliance (e.g.\ high-risk cases are always escalated) without
    filtering the LLM's predictions before they reach the log.
\end{enumerate}

The free-tier quota constraint on \texttt{gemini-2.0-flash} (20 requests/day)
directly shaped the fast-path design: by bypassing the LLM for 40--60\% of
production cases (those with high/critical risk or missing-input flags), the
daily quota is effectively reserved for genuinely ambiguous cases.

% ═════════════════════════════════════════════════════════════════════════════
\subsection{System Architecture}

\subsubsection{High-Level Design}

Every incoming \texttt{POST /predict} request is processed through one of two mutually
exclusive execution paths, with the decision made synchronously in \texttt{chain.py::\_fast\_path()}.

\begin{lstlisting}[style=ascii, caption={IAA execution flow}]
POST /predict
  |
  v
ProcessState validation (api.py or build_process_state from events)
  |
  v
run_iaa() [iaa.py]
  |
  +-- IAARetriever.retrieve() [retriever.py]
  |      embed query with all-MiniLM-L6-v2
  |      Qdrant top-k=2 cosine search
  |      returns list[dict] with id, document, metadata, distance
  |
  +-- run_chain() [chain.py]
        |
        +-- _fast_path(state, retrieved_ids)
              |
              +-- [FAST PATH] risk_severity in ("high","critical")
              |     -> escalate_to_human=True, model_used="fast_path/no_llm"
              |     -> return ActionRecommendation immediately (~0 ms)
              |
              +-- [FAST PATH] "missing"/"incomplete"/"required" in risk_flags
              |     -> action="request_missing_document", escalate=False
              |     -> return ActionRecommendation immediately (~0 ms)
              |
              +-- [STANDARD PATH] neither fast-path rule triggered
                    |
                    v
                 LLM.invoke(merged_prompt)
                    |
                    v
                 _extract_json(raw_response)
                    |
                    v
                 parse predicted_actions, task_list, workflow_timeline,
                 recommended_action  (all in one JSON response)
                    |
                    v
                 return ActionRecommendation (before overrides)
  |
  +-- apply_overrides(recommendation, state) [chain.py]
        |
        +-- HARD RULES (fire before confidence check):
        |     - high_risk_escalation: risk_severity in ("high","critical")
        |     - missing_critical_input: risk_flags contains missing/incomplete/required
        |
        +-- SOFT RULE: confidence < 0.45
        |     -> escalate_to_human=True, prediction preserved
        |
        +-- else: recommendation returned unchanged
  |
  v
_log(state, final) [iaa.py]
  append to data/predictions.jsonl
  |
  v
return ActionRecommendation
  |
  v
PredictResponse (api.py): {recommendation, latency_ms, corpus_size}
\end{lstlisting}

\subsubsection{Component Interaction}

The IAA codebase is decomposed into seven modules with strict calling relationships:

\begin{center}
\begin{tabular}{@{}lll@{}}
\toprule
\textbf{Module} & \textbf{Calls} & \textbf{Role} \\
\midrule
\texttt{api.py} & \texttt{iaa.py}, \texttt{schemas.py} & HTTP interface, request routing \\
\texttt{iaa.py} & \texttt{retriever.py}, \texttt{chain.py} & Core pipeline, logging, LangGraph node \\
\texttt{chain.py} & \texttt{schemas.py} & Fast path, LLM call, JSON parsing, overrides \\
\texttt{retriever.py} & \texttt{qdrant\_client\_singleton.py}, \texttt{schemas.py} & Vector search \\
\texttt{ingest.py} & \texttt{qdrant\_client\_singleton.py} & Corpus population \\
\texttt{schemas.py} & --- & Pydantic model definitions only \\
\texttt{qdrant\_client\_singleton.py} & --- & Shared Qdrant client + embedder \\
\bottomrule
\end{tabular}
\end{center}

Data flows in one direction: \texttt{api.py} receives an HTTP request, deserialises it
into \texttt{ProcessState} (directly or via \texttt{build\_process\_state} from raw events),
passes it to \texttt{run\_iaa()} in \texttt{iaa.py}. \texttt{iaa.py} calls
\texttt{IAARetriever.retrieve()} to get historical context, then calls
\texttt{run\_chain()} which either returns a fast-path result or invokes the LLM.
After \texttt{apply\_overrides()}, the final \texttt{ActionRecommendation} is logged
and returned up through \texttt{api.py} as a \texttt{PredictResponse}.

\subsubsection{Integration with Orchestrator}

The LangGraph orchestrator (\texttt{orchestrator/tools.py}) calls IAA via HTTP POST to
\texttt{http://localhost:8003/predict}. The payload is built by
\texttt{orchestrator/shims.py::build\_iaa\_input()}, which constructs a slim
\texttt{ProcessState} from AmpState fields: the last five case history events,
the last three email headers (subject/from/date only, no bodies), risk flags from RRA,
and governance decision from GPA. This bounded payload is the IAA latency fix --- it
prevents the LLM prompt from growing unboundedly with large case histories.

The raw HTTP response is unwrapped by \texttt{shims.py::iaa\_to\_amp()}, which extracts
the \texttt{recommendation} field from the \texttt{PredictResponse} wrapper and maps it
directly to \texttt{AmpState.iaa\_output}. The AmpState fields populated by IAA are:
\texttt{predicted\_actions}, \texttt{recommended\_action}, \texttt{override\_applied},
\texttt{override\_reason}, \texttt{retrieved\_context\_ids}, \texttt{escalate\_to\_human},
\texttt{operating\_context}, \texttt{model\_used}, \texttt{timestamp}, \texttt{task\_list},
and \texttt{workflow\_timeline}.

\textbf{Decision D1} (\texttt{context\_summary} removal): IAA's \texttt{ProcessState}
originally included a \texttt{context\_summary} field that was expected to be populated by
CSA. CSA never generated this field, causing a \texttt{KeyError} in the prompt assembly.
The field was removed from \texttt{ProcessState} as a required input and demoted to an
optional nullable field. The orchestrator shim was updated accordingly.

% ═════════════════════════════════════════════════════════════════════════════
\subsection{Database Design}

\subsubsection{Qdrant Vector Store}

IAA uses a single Qdrant collection named \texttt{iaa\_corpus}. All domains share this
collection; domain identity is encoded in payload metadata rather than separate collections.
This avoids collection sprawl and allows cross-domain retrieval when filtering is not applied.

\begin{center}
\begin{tabular}{@{}lll@{}}
\toprule
\textbf{Parameter} & \textbf{Value} & \textbf{Source} \\
\midrule
Collection name & \texttt{iaa\_corpus} & \texttt{ingest.py:COLLECTION} \\
Vector dimension & 384 & \texttt{ingest.py:VECTOR\_SIZE} \\
Distance metric & Cosine & \texttt{ingest.py::ensure\_collection()} \\
Embedding model & \texttt{all-MiniLM-L6-v2} & \texttt{qdrant\_client\_singleton.py} \\
Top-k at retrieval & 2 & \texttt{retriever.py:DEFAULT\_K} \\
Max doc chars in prompt & 400 & \texttt{retriever.py:MAX\_DOC\_CHARS} \\
\bottomrule
\end{tabular}
\end{center}

Each Qdrant point contains a vector (384-dim float) and a payload with the following fields:

\begin{center}
\begin{longtable}{@{}lll@{}}
\toprule
\textbf{Payload Field} & \textbf{Type} & \textbf{Description} \\
\midrule
\endhead
\texttt{document} & string & Natural-language process state description \\
\texttt{source} & string & Origin tag: \texttt{synthetic}, \texttt{bpic\_2012}, \texttt{bpic\_2018}, \texttt{cuad} \\
\texttt{domain} & string & Domain: \texttt{finance}, \texttt{legal}, \texttt{procurement}, \texttt{generic} \\
\texttt{process\_type} & string & Process name: \texttt{loan\_application}, \texttt{claim\_processing}, etc. \\
\texttt{next\_action} & string & Ground-truth next action observed in the historical log \\
\texttt{entity\_id} & string & Original case ID for audit traceability \\
\texttt{current\_event} & string & The current event type at the time of snapshot \\
\texttt{stage} & string & Process stage at time of snapshot \\
\texttt{actor} & string & Who performed the current event \\
\texttt{doc\_type} & string & \texttt{process\_trace}, \texttt{real\_trace}, \texttt{synthetic\_trace}, \texttt{policy\_clause} \\
\bottomrule
\end{longtable}
\end{center}

Documents are inserted via \texttt{ingest.py::\_batch\_add()}, which embeds documents
using \texttt{LocalEmbeddingFn} (wrapping \texttt{SentenceTransformer}), then upserts
batches of 100 \texttt{PointStruct} objects. Deduplication is two-stage: content-hash
UUIDs prevent in-batch duplicates, and a pre-flight \texttt{client.retrieve()} call
skips IDs already present in the collection.

\subsubsection{Singleton Client Design}

The embedded Qdrant file client takes an exclusive lock on the storage directory.
If two \texttt{QdrantClient(path=...)} instances are created in the same process
(e.g.\ one in \texttt{ingest.py} and one in \texttt{retriever.py}), the second raises:
\begin{quote}\ttfamily Storage folder ./qdrant\_store is already accessed by another instance\end{quote}

\texttt{qdrant\_client\_singleton.py::get\_qdrant\_client()} prevents this via a
module-level \texttt{\_client} global, initialised lazily on first call and reused
by all subsequent callers. The same pattern is applied to the embedding model
(\texttt{get\_embedding\_model()} with \texttt{\_embedding\_model} global), which
avoids the 90-second cold start cost of loading \texttt{all-MiniLM-L6-v2} more than once.

Connection mode is controlled by the \texttt{QDRANT\_HOST} environment variable:
if unset, the local file client is used at \texttt{QDRANT\_PATH} (default \texttt{./qdrant\_store});
if set (e.g.\ \texttt{qdrant} in Docker), a network client connects to port 6333.

\subsubsection{Corpus Sources}

\begin{center}
\begin{tabular}{@{}llll@{}}
\toprule
\textbf{Source} & \textbf{Format} & \textbf{Domain} & \textbf{Volume} \\
\midrule
Synthetic cases & JSON (200 cases) & claims, approvals, audits & \textasciitilde{}200 docs \\
BPIC 2012 & XES (loan applications) & finance & up to 500 cases \\
BPIC 2018 & XES (permit requests) & government & up to 500 cases \\
BPIC 2019 & XES (purchase orders) & procurement & up to 500 cases \\
CUAD & HuggingFace dataset & legal (contract clauses) & up to 200 docs \\
\bottomrule
\end{tabular}
\end{center}

Synthetic cases are JSON files at \texttt{data/synthetic\_cases.json} (200 cases,
generated by \texttt{scripts/generate\_synthetic.py}). Each case contains:
\texttt{entity\_id}, \texttt{event\_type}, \texttt{prior\_actions}, \texttt{next\_action},
\texttt{stage}, \texttt{risk\_severity}, \texttt{process\_type}, \texttt{domain},
and \texttt{actor\_sequence} (list of \texttt{\{event\_type, actor\}} dicts).

The audit log is written to \texttt{data/predictions.jsonl} (configured via the
\texttt{LOG\_FILE} env var). Each line is a JSON object written by \texttt{iaa.py::\_log()}
containing: \texttt{timestamp}, \texttt{entity\_id}, \texttt{current\_stage},
\texttt{current\_event}, \texttt{operating\_context}, \texttt{recommended},
\texttt{confidence}, \texttt{override\_applied}, \texttt{override\_reason},
\texttt{escalate}, \texttt{model\_used}, \texttt{task\_count}, \texttt{timeline\_steps}.

% ═════════════════════════════════════════════════════════════════════════════
\subsection{Code File Explanations}

\subsubsection{\texttt{api.py}}

\textbf{Purpose.} FastAPI REST server providing the HTTP interface consumed by the
LangGraph orchestrator. Runs on port 8003 in the AMP stack.

\textbf{Key functions/classes:}
\begin{itemize}[nosep]
  \item \texttt{PredictRequest} --- accepts either a pre-built \texttt{ProcessState}
    or a raw list of \texttt{EventLogEntry} objects; optional \texttt{domain\_filter}
    and \texttt{process\_type\_filter} for scoped retrieval.
  \item \texttt{PredictResponse} --- wraps \texttt{ActionRecommendation} with
    \texttt{latency\_ms} and optional \texttt{corpus\_size}.
  \item \texttt{predict(req: PredictRequest) -> PredictResponse} --- main endpoint;
    builds \texttt{ProcessState} from events if needed, calls \texttt{run\_iaa()},
    measures wall-clock latency, appends corpus count.
  \item \texttt{ingest\_documents(req: IngestRequest, background\_tasks)} --- accepts
    raw documents + metadata, kicks off \texttt{\_batch\_add()} asynchronously so
    the HTTP response is immediate.
  \item \texttt{startup()} --- pre-warms \texttt{IAARetriever} and LLM on FastAPI
    startup, reducing first-call latency.
  \item \texttt{health()} --- liveness check returning \texttt{status: ok} and
    the active \texttt{LLM\_PROVIDER}.
\end{itemize}

\textbf{Notable details.} The \texttt{/workflow} endpoint delegates to
\texttt{orchestrator.orchestrator.run\_workflow()} (the IAA's own local orchestrator,
not the AMP LangGraph orchestrator). This endpoint is used by the CLI demo but not
by the production AMP pipeline, which calls \texttt{/predict} directly.

\subsubsection{\texttt{chain.py}}

\textbf{Purpose.} Implements the LLM chain: a single merged prompt replacing the
original three-step chain, plus the fast-path bypass and post-LLM rule overrides.

\textbf{Key functions/classes:}
\begin{itemize}[nosep]
  \item \texttt{get\_llm(temperature=0.1)} --- provider factory; returns
    \texttt{ChatGoogleGenerativeAI} or \texttt{ChatOpenAI} based on \texttt{LLM\_PROVIDER}.
    Temperature fixed at 0.1 for determinism.
  \item \texttt{\_fast\_path(state, retrieved\_ids) -> ActionRecommendation | None} ---
    returns a fully-formed \texttt{ActionRecommendation} with
    \texttt{model\_used="fast\_path/no\_llm"} when either hard rule fires; returns
    \texttt{None} if standard path is required.
  \item \texttt{run\_chain(state, retrieved\_context, retrieved\_ids, llm)} ---
    checks fast path first; on standard path, formats the merged prompt
    (\texttt{MERGED\_SYSTEM} + \texttt{MERGED\_USER}), calls \texttt{llm.invoke()},
    extracts JSON via \texttt{\_extract\_json()}, and parses
    \texttt{predicted\_actions}, \texttt{task\_list}, \texttt{workflow\_timeline},
    and \texttt{recommended\_action}.
  \item \texttt{apply\_overrides(recommendation, state) -> ActionRecommendation} ---
    iterates \texttt{HARD\_RULES}, then checks confidence threshold (0.45).
    Returns a new immutable \texttt{ActionRecommendation}; never mutates input.
  \item \texttt{\_extract\_json(text)} --- strips markdown fences (\texttt{```json}),
    uses regex to find the first \texttt{\{...\}} block, and raises \texttt{ValueError}
    if no JSON is found.
\end{itemize}

\textbf{HARD\_RULES} list contains three entries:
\begin{enumerate}[nosep]
  \item \texttt{high\_risk\_escalation}: fires when \texttt{risk\_severity} is
    \texttt{"high"} or \texttt{"critical"}, action = \texttt{escalate\_to\_human},
    confidence = 0.99, escalate = True.
  \item \texttt{missing\_critical\_input}: fires when any word in
    \texttt{risk\_flags} contains \texttt{"missing"}, \texttt{"incomplete"}, or
    \texttt{"required"}, action = \texttt{request\_missing\_document},
    confidence = 0.95, escalate = False.
  \item \texttt{low\_confidence\_escalation}: not evaluated via the condition lambda
    (lambda always returns \texttt{False}); the confidence check is handled separately
    as a soft escalation --- the original prediction is preserved.
\end{enumerate}

\textbf{CONFIDENCE\_THRESHOLD} is 0.45; predictions below this set
\texttt{escalate\_to\_human=True} without replacing the recommended action.

\textbf{Notable quirk.} The merged prompt instructs the LLM to ``return valid JSON
only, no markdown'', but in practice Gemini occasionally wraps output in
\texttt{```json} fences. \texttt{\_extract\_json()} handles this defensively.

\subsubsection{\texttt{retriever.py}}

\textbf{Purpose.} Domain-agnostic RAG retrieval: embeds the current process state
as a natural-language query, searches Qdrant for the most similar historical
process snapshots, and formats results for LLM prompt injection.

\textbf{Key functions/classes:}
\begin{itemize}[nosep]
  \item \texttt{build\_query(state: ProcessState) -> str} --- constructs a plain-text
    search query from current stage, current event, last three prior actions, risk
    severity, risk flags, and optional context summary.
  \item \texttt{IAARetriever} class:
    \begin{itemize}[nosep]
      \item \texttt{retrieve(state, domain\_filter=None, process\_type\_filter=None) -> list[dict]} ---
        embeds the query with \texttt{get\_embedding\_model()}, calls
        \texttt{client.query\_points()} with \texttt{limit=k}, applies optional
        payload filters. Graceful degradation: returns \texttt{[]} on any exception.
      \item \texttt{format\_for\_prompt(retrieved: list[dict]) -> str} --- formats
        each result as a numbered \texttt{[Context N]} block with source, process type,
        next-action label, and truncated document body (\texttt{MAX\_DOC\_CHARS=400}).
    \end{itemize}
\end{itemize}

\textbf{Notable details.} \texttt{DEFAULT\_K=2} (reduced from 4 in the Phase 1
latency fix) keeps the injected context small. The retriever client is lazy-loaded:
the property \texttt{IAARetriever.client} initialises the singleton and calls
\texttt{ensure\_collection()} on first access.

\subsubsection{\texttt{iaa.py}}

\textbf{Purpose.} Core IAA entrypoint. Exposes two interfaces:
\texttt{run\_iaa()} for direct invocation and \texttt{iaa\_node()} as a
LangGraph-compatible node.

\textbf{Key functions:}
\begin{itemize}[nosep]
  \item \texttt{build\_process\_state(events, context\_summary=None, risk\_flags=None, ...)} ---
    sorts \texttt{EventLogEntry} objects by timestamp, treats the last event as the
    current state and all prior events as history. Infers \texttt{operating\_context}:
    \texttt{"full"} only when risk severity, policy constraints, and assigned owner are
    all provided; otherwise \texttt{"partial --- <missing agents>"}.
  \item \texttt{run\_iaa(state, domain\_filter=None, process\_type\_filter=None, retriever=None, llm=None)} ---
    lazy-initialises the module-level \texttt{\_retriever} and \texttt{\_llm} singletons,
    retrieves context, runs the chain, applies overrides, logs the result.
  \item \texttt{iaa\_node(graph\_state: dict) -> dict} --- reads
    \texttt{graph\_state["process\_state"]}, optionally supplements it with
    \texttt{csa\_output["context\_summary"]} and \texttt{rra\_output["severity"]},
    calls \texttt{run\_iaa()}, writes \texttt{iaa\_recommendation},
    \texttt{iaa\_routing\_action}, and \texttt{iaa\_escalate} to the state dict.
    Returns a fallback escalation if \texttt{process\_state} is absent.
  \item \texttt{\_log(state, rec)} --- appends a JSON line to \texttt{LOG\_FILE}
    (default \texttt{./data/predictions.jsonl}), creating the directory if needed.
    Failures are suppressed with a warning log.
\end{itemize}

\textbf{Notable detail.} Module-level \texttt{\_retriever} and \texttt{\_llm} are
initialised to \texttt{None} and populated on first \texttt{run\_iaa()} call.
Injecting these via the function parameters (as tests do) bypasses the globals,
enabling test isolation without patching module state.

\subsubsection{\texttt{schemas.py}}

\textbf{Purpose.} All Pydantic v2 models for IAA input and output. No business logic.

\textbf{Key classes:}
\begin{itemize}[nosep]
  \item \texttt{EventLogEntry} --- universal event log row; accepts BPIC XES data,
    CUAD events, or custom logs. Required fields: \texttt{event\_id}, \texttt{entity\_id},
    \texttt{event\_type}, \texttt{stage}, \texttt{timestamp}.
  \item \texttt{ProcessStep} --- structured prior step with \texttt{event\_type},
    \texttt{actor}, \texttt{timestamp}, \texttt{outcome}.
  \item \texttt{ProcessState} --- primary input to the chain; assembled from event
    history. Contains optional upstream agent fields (\texttt{risk\_flags},
    \texttt{risk\_severity}, \texttt{policy\_constraints}, \texttt{assigned\_owner}).
  \item \texttt{TaskItem} --- one task for DA: \texttt{task\_id}, \texttt{action},
    \texttt{assigned\_role}, \texttt{priority}, \texttt{depends\_on}, \texttt{reasoning}.
  \item \texttt{WorkflowStep} --- one step in the timeline: \texttt{step},
    \texttt{event\_type}, \texttt{actor}, \texttt{status}, \texttt{timestamp}.
  \item \texttt{PredictedAction} --- one ranked action: \texttt{action\_id},
    \texttt{label}, \texttt{confidence} (0.0--1.0 constrained), \texttt{reasoning}.
  \item \texttt{ActionRecommendation} --- top-level output; \texttt{predicted\_actions}
    is constrained \texttt{min\_length=1, max\_length=3}.
  \item \texttt{IAAState} --- LangGraph shared state contract; serves as the integration
    contract for the AITS team.
\end{itemize}

\subsubsection{\texttt{ingest.py}}

\textbf{Purpose.} Universal corpus ingestion pipeline. Converts BPIC XES files,
CUAD clauses, tabular CSV/JSON logs, and synthetic cases into Qdrant points.

\textbf{Key functions:}
\begin{itemize}[nosep]
  \item \texttt{ensure\_collection(client)} --- idempotent collection creation.
  \item \texttt{\_make\_id(text, prefix) -> str} --- deterministic MD5-based UUID
    from content + prefix; identical content with different prefixes yields different IDs.
  \item \texttt{ingest\_bpic\_xes(xes\_path, source\_name, process\_type, domain, max\_cases, client) -> int} ---
    reads XES via pm4py, converts to DataFrame, samples up to 500 cases, generates
    process-state snapshots (one per event step, up to 8 per case), embeds and upserts.
  \item \texttt{ingest\_cuad(client=None, max\_docs=200) -> int} --- loads CUAD from
    HuggingFace, wraps each contract clause as a process document.
  \item \texttt{ingest\_tabular(file\_path, column\_map, source\_name, domain, process\_type, ...) -> int} ---
    generic entry point for client-specific CSV or JSON event logs via a column mapping dict.
  \item \texttt{ingest\_synthetic\_cases(cases, source\_name, domain, process\_type, client) -> int} ---
    ingests pre-built synthetic case dicts. Includes \texttt{entity\_id} in the hash
    to prevent dedup collisions on cases with identical event sequences.
  \item \texttt{bootstrap\_corpus(include\_synthetic, include\_bpic, include\_cuad, ...)} ---
    one-shot setup function; fastest path is \texttt{bootstrap\_corpus(include\_synthetic=True)}.
  \item \texttt{\_batch\_add(client, docs, ids, metas, batch\_size=100)} --- two-stage
    dedup then embed + upsert in batches of 100.
\end{itemize}

\subsubsection{\texttt{qdrant\_client\_singleton.py}}

\textbf{Purpose.} Process-wide singleton for the Qdrant client and embedding model,
preventing file lock conflicts in local mode.

\textbf{Key functions/classes:}
\begin{itemize}[nosep]
  \item \texttt{get\_qdrant\_client() -> QdrantClient} --- returns the process-wide
    client; creates it on first call. Uses network client if \texttt{QDRANT\_HOST}
    is set, otherwise file client at \texttt{QDRANT\_PATH}.
  \item \texttt{get\_embedding\_model() -> SentenceTransformer} --- lazy-loads
    \texttt{all-MiniLM-L6-v2}; cached in \texttt{\_embedding\_model} global.
  \item \texttt{LocalEmbeddingFn} --- wrapper class with \texttt{\_\_call\_\_},
    \texttt{embed\_documents()}, and \texttt{embed\_query()} so it works as both
    a callable and a LangChain-compatible embedder.
\end{itemize}

\textbf{Config constants:} \texttt{QDRANT\_PATH} (env: \texttt{QDRANT\_PATH},
default \texttt{./qdrant\_store}), \texttt{QDRANT\_HOST} (env: \texttt{QDRANT\_HOST},
default \texttt{None}), \texttt{QDRANT\_PORT} (env: \texttt{QDRANT\_PORT}, default 6333).

\subsubsection{\texttt{conftest.py}}

\textbf{Purpose.} pytest configuration that mocks the LLM globally so all 77 unit
tests run without any API key.

\textbf{Mechanism.} The \texttt{mock\_llm\_globally} fixture is marked
\texttt{autouse=True}, applying it to every test automatically. Tests marked
\texttt{@pytest.mark.integration} receive a \texttt{yield} without any mock applied,
allowing real API calls. Non-integration tests receive a \texttt{MagicMock} LLM whose
\texttt{invoke.side\_effect} cycles through three pre-baked responses:

\begin{enumerate}[nosep]
  \item \texttt{MOCK\_STEP1} --- plain-text process summary string.
  \item \texttt{MOCK\_STEP2} --- JSON with \texttt{predicted\_actions},
    \texttt{task\_list}, \texttt{workflow\_timeline}.
  \item \texttt{MOCK\_STEP3} --- JSON with \texttt{action\_id}, \texttt{label},
    \texttt{confidence}, \texttt{justification}.
\end{enumerate}

The patch target is \texttt{agent.chain.get\_llm}; since \texttt{run\_chain()} calls
\texttt{get\_llm()} only when \texttt{llm} is not injected, tests that pass a mock
LLM directly bypass the patch entirely.

% ═════════════════════════════════════════════════════════════════════════════
\subsection{Input Specification}

\subsubsection{PredictRequest Schema}

\begin{longtable}{@{}p{3.8cm}p{3cm}p{1.5cm}p{6cm}@{}}
\toprule
\textbf{Field} & \textbf{Type} & \textbf{Required} & \textbf{Description} \\
\midrule
\endhead
\texttt{process\_state} & \texttt{ProcessState} & No* & Pre-built process state (path 1) \\
\texttt{events} & \texttt{list[EventLogEntry]} & No* & Raw event log (path 2 --- IAA builds state) \\
\texttt{context\_summary} & \texttt{str | None} & No & Optional CSA context summary \\
\texttt{risk\_flags} & \texttt{list[str] | None} & No & RRA risk flags \\
\texttt{risk\_severity} & \texttt{str | None} & No & RRA severity: low/medium/high/critical \\
\texttt{policy\_constraints} & \texttt{list[str] | None} & No & GPA active constraints \\
\texttt{assigned\_owner} & \texttt{str | None} & No & DA assigned owner (if pre-assigned) \\
\texttt{domain\_filter} & \texttt{str | None} & No & Scope Qdrant search to this domain \\
\texttt{process\_type\_filter} & \texttt{str | None} & No & Scope Qdrant search to this process type \\
\bottomrule
\multicolumn{4}{@{}l@{}}{\small * Either \texttt{process\_state} or \texttt{events} must be provided; both absent returns HTTP 422.}
\end{longtable}

\subsubsection{ProcessState Schema}

\begin{longtable}{@{}p{3.8cm}p{3.2cm}p{1.5cm}p{5.5cm}@{}}
\toprule
\textbf{Field} & \textbf{Type} & \textbf{Required} & \textbf{Description} \\
\midrule
\endhead
\texttt{entity\_id} & \texttt{str} & Yes & Case/claim/contract identifier \\
\texttt{current\_stage} & \texttt{str} & Yes & Current process stage \\
\texttt{current\_event\_type} & \texttt{str} & Yes & Most recent event type \\
\texttt{prior\_actions} & \texttt{list[str]} & No & Ordered prior event types (default: \texttt{[]}) \\
\texttt{prior\_steps} & \texttt{list[ProcessStep]} & No & Structured steps with actor/timestamp \\
\texttt{actor} & \texttt{str | None} & No & Who performed the current event \\
\texttt{timestamp} & \texttt{datetime} & Yes & When the current event occurred \\
\texttt{metadata} & \texttt{dict} & No & Domain-specific extras (amounts, priority, etc.) \\
\texttt{context\_summary} & \texttt{str | None} & No & CSA summary (optional, rarely populated) \\
\texttt{risk\_flags} & \texttt{list[str]} & No & RRA flags (default: \texttt{[]}) \\
\texttt{risk\_severity} & \texttt{str | None} & No & \texttt{low | medium | high | critical} \\
\texttt{policy\_constraints} & \texttt{list[str]} & No & GPA constraints (default: \texttt{[]}) \\
\texttt{assigned\_owner} & \texttt{str | None} & No & Pre-assigned owner (optional) \\
\texttt{operating\_context} & \texttt{str} & No & \texttt{"full"} or \texttt{"partial --- <agents>"} \\
\bottomrule
\end{longtable}

The \texttt{risk\_severity} field directly controls fast-path routing. Values
\texttt{"high"} and \texttt{"critical"} trigger the \texttt{high\_risk\_escalation}
fast-path rule, bypassing the LLM entirely and returning
\texttt{escalate\_to\_human=True} with confidence 0.99. Values \texttt{"low"} and
\texttt{"medium"} take the standard path. A \texttt{None} value also takes the
standard path but the LLM receives \texttt{"unknown"} for the risk level field in
the prompt, potentially reducing prediction quality.

The \texttt{risk\_flags} field triggers the second fast-path rule when any flag
contains the substrings \texttt{"missing"}, \texttt{"incomplete"}, or
\texttt{"required"} (case-insensitive substring match on the joined flags string).
This covers common flag patterns such as \texttt{"missing\_required\_attachment"},
\texttt{"incomplete\_documentation"}, and \texttt{"required\_field\_absent"}.

\subsubsection{Example Inputs}

\begin{lstlisting}[caption={Standard path input: medium-risk vendor onboarding}]
{
  "process_state": {
    "entity_id":          "CASE-VND-001",
    "current_stage":      "document_verification",
    "current_event_type": "document_submitted",
    "prior_actions":      ["vendor_profile_created", "initial_review"],
    "actor":              "procurement_officer",
    "timestamp":          "2026-04-30T10:00:00Z",
    "metadata":           {"monetary_value": 150000},
    "risk_flags":         ["new_vendor", "high_value"],
    "risk_severity":      "medium",
    "policy_constraints": [],
    "operating_context":  "full"
  },
  "domain_filter":       "procurement",
  "process_type_filter": "vendor_onboarding"
}
\end{lstlisting}

\begin{lstlisting}[caption={Fast-path input: critical-risk fraud investigation}]
{
  "events": [
    {
      "event_id":   "e1",
      "entity_id":  "CASE-FRD-003",
      "event_type": "suspicious_transaction_flagged",
      "stage":      "fraud_detection",
      "timestamp":  "2026-04-30T09:00:00Z",
      "metadata":   {"amount": 12000}
    },
    {
      "event_id":   "e2",
      "entity_id":  "CASE-FRD-003",
      "event_type": "manual_review_initiated",
      "stage":      "investigation",
      "timestamp":  "2026-04-30T09:30:00Z",
      "metadata":   {}
    }
  ],
  "risk_flags":    ["fraud_risk", "regulatory_risk"],
  "risk_severity": "critical"
}
\end{lstlisting}

% ═════════════════════════════════════════════════════════════════════════════
\subsection{Output Specification}

\subsubsection{PredictResponse Schema}

\begin{center}
\begin{tabular}{@{}llll@{}}
\toprule
\textbf{Field} & \textbf{Type} & \textbf{Nullable} & \textbf{Description} \\
\midrule
\texttt{recommendation} & \texttt{ActionRecommendation} & No & Full IAA output \\
\texttt{latency\_ms} & \texttt{float} & No & Wall-clock time for the full request \\
\texttt{corpus\_size} & \texttt{int | None} & Yes & Number of docs in Qdrant at time of call \\
\bottomrule
\end{tabular}
\end{center}

\subsubsection{ActionRecommendation Schema}

\begin{longtable}{@{}p{4.2cm}p{3.5cm}p{6cm}@{}}
\toprule
\textbf{Field} & \textbf{Type} & \textbf{Description} \\
\midrule
\endhead
\texttt{predicted\_actions} & \texttt{list[PredictedAction]} & Top 1--3 actions ranked by confidence descending \\
\texttt{recommended\_action} & \texttt{PredictedAction} & The single best action to surface \\
\texttt{override\_applied} & \texttt{bool} & True if a hard rule fired (default: False) \\
\texttt{override\_reason} & \texttt{str | None} & Human-readable rule name + reason \\
\texttt{retrieved\_context\_ids} & \texttt{list[str]} & Qdrant point IDs used for retrieval \\
\texttt{escalate\_to\_human} & \texttt{bool} & True if human review required (default: False) \\
\texttt{operating\_context} & \texttt{str} & Passed from ProcessState for audit \\
\texttt{model\_used} & \texttt{str} & e.g.\ \texttt{gemini/gemini-2.0-flash} or \texttt{fast\_path/no\_llm} \\
\texttt{timestamp} & \texttt{datetime} & When this recommendation was generated \\
\texttt{task\_list} & \texttt{list[TaskItem]} & Tasks for Delegation Agent to consume \\
\texttt{workflow\_timeline} & \texttt{list[WorkflowStep]} & Chronological case timeline \\
\bottomrule
\end{longtable}

Confidence scoring uses the raw LLM output's confidence float (0.0--1.0), constrained
by the \texttt{PredictedAction} field validator. The LLM is instructed to rank
\texttt{predicted\_actions} descending by confidence and to include
\texttt{"escalate\_for\_review"} as the last action when any confidence falls below 0.3.
The \texttt{escalate\_to\_human} flag on \texttt{ActionRecommendation} is set to True
when: (a) the recommended action's \texttt{action\_id} is \texttt{"escalate\_for\_review"}
or \texttt{"escalate\_to\_human"}, (b) the fast-path high-risk rule fires, or (c) the
confidence threshold override fires.

\subsubsection{WorkflowTimeline and TaskList}

\textbf{WorkflowStep} fields:
\begin{center}
\begin{tabular}{@{}llll@{}}
\toprule
\textbf{Field} & \textbf{Type} & \textbf{Required} & \textbf{Description} \\
\midrule
\texttt{step} & \texttt{int} & Yes & Ordinal position (1-indexed) \\
\texttt{event\_type} & \texttt{str} & Yes & Event type label \\
\texttt{actor} & \texttt{str | None} & No & Role responsible for this step \\
\texttt{status} & \texttt{str} & Yes & \texttt{completed | current | predicted} \\
\texttt{timestamp} & \texttt{datetime | None} & No & Null for predicted steps \\
\bottomrule
\end{tabular}
\end{center}

\textbf{TaskItem} fields:
\begin{center}
\begin{tabular}{@{}llll@{}}
\toprule
\textbf{Field} & \textbf{Type} & \textbf{Required} & \textbf{Description} \\
\midrule
\texttt{task\_id} & \texttt{str} & Yes & Unique ID e.g.\ \texttt{T001} \\
\texttt{action} & \texttt{str} & Yes & Snake\_case action identifier \\
\texttt{assigned\_role} & \texttt{str | None} & No & Role for DA to consume \\
\texttt{priority} & \texttt{str} & No & \texttt{high | medium | low} (default: \texttt{medium}) \\
\texttt{depends\_on} & \texttt{list[str]} & No & task\_ids this task depends on \\
\texttt{reasoning} & \texttt{str} & Yes & Why this task is included \\
\bottomrule
\end{tabular}
\end{center}

The Streamlit dashboard (\texttt{demo/app.py}) renders the \texttt{workflow\_timeline}
as a status list with icons (\texttt{\checkmark} for completed, \texttt{▶} for current,
\texttt{○} for predicted) and renders \texttt{task\_list} as a \texttt{st.dataframe}
table in the IAA tab. The DA uses the first \texttt{TaskItem.assigned\_role} from
\texttt{task\_list[0]} to populate \texttt{recommended\_candidates} in its input.

\subsubsection{Fast-Path vs Standard-Path Output Differences}

\begin{center}
\begin{tabular}{@{}lll@{}}
\toprule
\textbf{Field} & \textbf{Fast Path} & \textbf{Standard Path} \\
\midrule
\texttt{predicted\_actions} & Exactly 1 item & 1--3 items \\
\texttt{task\_list} & \texttt{[]} (empty) & Populated by LLM \\
\texttt{workflow\_timeline} & \texttt{[]} (empty) & Populated by LLM \\
\texttt{override\_applied} & \texttt{True} & \texttt{False} (unless soft rule fires) \\
\texttt{model\_used} & \texttt{"fast\_path/no\_llm"} & \texttt{"gemini/gemini-2.0-flash"} \\
\texttt{latency\_ms} & \textasciitilde{}0--5\,ms & \textasciitilde{}3000--10000\,ms \\
\texttt{escalate\_to\_human} & \texttt{True} (risk) or \texttt{False} (missing doc) & Depends on LLM + overrides \\
\bottomrule
\end{tabular}
\end{center}

% ═════════════════════════════════════════════════════════════════════════════
\subsection{Human-in-the-Loop Conditions}

The \texttt{escalate\_to\_human} flag on \texttt{ActionRecommendation} is set to
\texttt{True} under the following conditions, in priority order:

\begin{enumerate}
  \item \textbf{Fast-path rule 1 --- high/critical risk} (\texttt{chain.py::\_fast\_path()}).
    When \texttt{ProcessState.risk\_severity} is \texttt{"high"} or \texttt{"critical"},
    the LLM is bypassed entirely. \texttt{escalate\_to\_human=True},
    \texttt{recommended\_action.action\_id="escalate\_to\_human"},
    \texttt{confidence=0.99}, \texttt{model\_used="fast\_path/no\_llm"}.

  \item \textbf{Fast-path rule 2 --- missing required inputs} (\texttt{chain.py::\_fast\_path()}).
    When any string in \texttt{risk\_flags} contains \texttt{"missing"},
    \texttt{"incomplete"}, or \texttt{"required"} (case-insensitive substring match
    on the space-joined flags). The LLM is bypassed.
    \texttt{escalate\_to\_human=False}, \texttt{recommended\_action.action\_id="request\_missing\_document"}.
    Note: this fast-path rule does \emph{not} set escalate; it routes to document
    retrieval instead. However it preserves the same output shape.

  \item \textbf{Post-LLM hard rule 1 --- high/critical risk override}
    (\texttt{chain.py::apply\_overrides()}).
    If the fast path was not triggered (e.g.\ the fast path check ran but
    risk\_severity changed between the fast-path check and override application,
    which should not normally occur), \texttt{apply\_overrides()} re-applies the
    same \texttt{high\_risk\_escalation} rule. In practice, fast path always fires first
    for high/critical risk, so this override is a belt-and-suspenders safety net.

  \item \textbf{Post-LLM hard rule 2 --- missing input override}
    (\texttt{chain.py::apply\_overrides()}).
    Same condition as fast-path rule 2, applied post-LLM in the rare case where
    fast path did not fire but the state has missing-input flags.

  \item \textbf{Post-LLM soft rule --- low confidence} (\texttt{chain.py::apply\_overrides()}).
    When \texttt{recommended\_action.confidence < CONFIDENCE\_THRESHOLD} (0.45),
    \texttt{escalate\_to\_human=True} is set and \texttt{override\_applied=True},
    but the original \texttt{recommended\_action} is \emph{preserved} (not replaced).
    The override reason records the exact confidence and threshold values.

  \item \textbf{Standard-path auto-escalation via action\_id}
    (\texttt{chain.py::run\_chain()} line 298).
    When the LLM's \texttt{recommended\_action.action\_id} is either
    \texttt{"escalate\_for\_review"} or \texttt{"escalate\_to\_human"},
    \texttt{escalate\_to\_human=True} is set on the returned \texttt{ActionRecommendation}
    without any override rule firing. The LLM itself chose to escalate.
\end{enumerate}

\textbf{Downstream effect.} When the orchestrator receives
\texttt{escalate\_to\_human=True} in the IAA output, the \texttt{iaa\_to\_amp()} shim
propagates it to \texttt{AmpState.iaa\_output.escalate\_to\_human}. Decision D7
then sets \texttt{AmpState.meta.requires\_escalation=True} if any of IAA, RRA, GPA,
or DA has an escalation signal. The Delegation Agent receives this flag and sets
\texttt{da\_output.escalation\_required=True}, routing the case to the
\texttt{human\_owner} channel rather than an automated agent or queue.

% ═════════════════════════════════════════════════════════════════════════════
\subsection{Role in the Multi-Agent System}

IAA is the second stage of the AMP pipeline (Stage 2), running concurrently with
RRA and GPA after CSA completes. It does \emph{not} call CSA, RRA, or GPA directly;
it consumes pre-populated fields from \texttt{AmpState} via orchestrator shims.

\begin{center}
\begin{tabular}{@{}lll@{}}
\toprule
\textbf{Agent} & \textbf{Direction} & \textbf{What IAA gets/gives} \\
\midrule
CSA (Stage 1) & CSA $\to$ IAA & Last 5 case history events + last 3 email headers \\
RRA (Stage 2, parallel) & RRA $\to$ IAA & \texttt{risk\_flags}, \texttt{risk\_severity} \\
GPA (Stage 2, parallel) & GPA $\to$ IAA & \texttt{governance\_decision}, \texttt{requires\_approval} \\
DA (Stage 3) & IAA $\to$ DA & \texttt{recommended\_action}, \texttt{task\_list}, \texttt{escalate\_to\_human} \\
\bottomrule
\end{tabular}
\end{center}

The orchestrator's \texttt{build\_iaa\_input()} shim constructs a slim \texttt{ProcessState}
from AmpState: it truncates case history to the last 5 events (structured as
\texttt{\{timestamp, actor, event\_type, status\}}) and emails to the last 3 headers
only (no bodies, no attachments). This is the IAA latency fix --- it prevents the
prompt from growing with large case histories.

RRA and GPA run in parallel with IAA via \texttt{asyncio.gather()} in the orchestrator.
IAA therefore does not have access to RRA or GPA results at call time. The orchestrator
may provide \texttt{risk\_flags}/\texttt{risk\_severity} from a fast-completing RRA
call, but in the production flow these are typically absent or null when IAA starts,
so IAA infers \texttt{operating\_context="partial"} and proceeds without them.

DA consumes IAA's \texttt{recommended\_action.action\_id} (mapped via
\texttt{build\_da\_input()} shim as \texttt{iaa\_recommended\_action}),
the \texttt{assigned\_role} from \texttt{task\_list[0]} (mapped as
\texttt{iaa\_recommended\_action.assigned\_role}), and the \texttt{escalate\_to\_human}
flag (mapped into \texttt{escalate\_flags.iaa}). The DA uses these three signals,
along with RRA's \texttt{risk\_level} and GPA's \texttt{governance\_decision}, to
make its final routing decision.

% ═════════════════════════════════════════════════════════════════════════════
\subsection{Key Features}

\subsubsection{Fast-Path Escalation}

The fast path (\texttt{chain.py::\_fast\_path()}) checks two conditions before any
embedding or LLM call: whether \texttt{risk\_severity} is \texttt{"high"} or
\texttt{"critical"}, and whether any \texttt{risk\_flag} contains keywords suggesting
missing inputs. For the high-risk case, it returns an \texttt{ActionRecommendation}
with \texttt{action\_id="escalate\_to\_human"}, \texttt{confidence=0.99}, and
\texttt{model\_used="fast\_path/no\_llm"} in approximately 0--5\,ms wall-clock time,
compared to 3,000--10,000\,ms for a standard Gemini call. The business rationale is
twofold: first, high-risk cases should always go to humans regardless of what the
LLM might predict, so the LLM call would be wasted computation; second, on the Gemini
free tier (20 requests/day), preserving quota for genuinely ambiguous medium-risk
cases is operationally critical. The fast path was introduced in Session 3 of the
project to address quota exhaustion during demo runs with multiple fraud investigation
cases.

\subsubsection{RAG-Augmented Prediction}

The retrieval step (\texttt{retriever.py::IAARetriever.retrieve()}) embeds the current
process state as a natural-language query using \texttt{all-MiniLM-L6-v2} (384-dim,
local, no API calls) and searches the \texttt{iaa\_corpus} Qdrant collection for the
top-$k$ most similar historical process snapshots (default $k=2$, reduced from 4 in
the Phase 1 latency fix). Each retrieved document contains the process type, current
stage, prior action sequence, actor, and the observed next action --- exactly the
ground-truth signal needed to ground the LLM's prediction in real workflow patterns.
Retrieved documents are formatted by \texttt{format\_for\_prompt()} as numbered
\texttt{[Context N]} blocks with their metadata and truncated body (400 chars max),
then injected into the \texttt{MERGED\_USER} prompt. The retriever degrades gracefully:
if the corpus is empty or the Qdrant search fails, it returns an empty list and
the LLM proceeds with no historical context.

\subsubsection{Rule-Based Post-LLM Overrides}

The override layer (\texttt{chain.py::apply\_overrides()}) runs after the LLM
chain returns, not before. This design choice is deliberate: running overrides
after the LLM ensures that the original LLM predictions are always logged in
\texttt{predicted\_actions}, enabling post-hoc analysis of model behaviour even
for overridden cases. The two hard rules (high risk and missing inputs) \emph{replace}
the \texttt{recommended\_action} entirely while preserving the LLM's \texttt{predicted\_actions}
in the returned object. The soft rule (low confidence) sets the \texttt{escalate\_to\_human}
flag without replacing the prediction. The \texttt{task\_list} and
\texttt{workflow\_timeline} are explicitly threaded through all override paths to
prevent their loss.

\subsubsection{Structured Output for Downstream Consumption}

The merged LLM prompt (\texttt{MERGED\_USER}) instructs the model to produce
\texttt{task\_list} and \texttt{workflow\_timeline} alongside the action prediction,
in the same JSON response. The \texttt{task\_list} gives DA a structured list of
tasks with \texttt{assigned\_role} values, enabling intelligent routing beyond simple
escalation. The \texttt{workflow\_timeline} provides the Streamlit dashboard with a
chronological view of the case from start to end, with each step labelled as
\texttt{completed}, \texttt{current}, or \texttt{predicted}. Both structures default
to empty lists on the fast path and when the LLM omits them.

\subsubsection{Configurable LLM Backend}

The LLM provider is controlled by the \texttt{LLM\_PROVIDER} environment variable
(\texttt{"gemini"} or \texttt{"openai"}). The factory function \texttt{get\_llm()}
in \texttt{chain.py} instantiates the appropriate LangChain chat model. The Gemini
model name is further configurable via \texttt{GEMINI\_MODEL} (default
\texttt{"gemini-2.0-flash"}), enabling model switching without code changes. During
Session 3, the default was changed from \texttt{gemini-2.5-flash} to
\texttt{gemini-2.0-flash} as the 2.0 model had more available quota on the free tier.
The OpenAI path instantiates \texttt{gpt-4o-mini} but has not been tested end-to-end;
it is present as a fallback for environments without Gemini access.

\subsubsection{Test Isolation}

The \texttt{conftest.py} \texttt{autouse} fixture patches \texttt{agent.chain.get\_llm}
for every non-integration test, returning a \texttt{MagicMock} whose \texttt{invoke}
cycles through three pre-baked JSON strings. This design means no network connection,
no API key, and no rate limits are required for the 77-test unit suite. The retriever
is not mocked at the conftest level --- tests that need a mock retriever inject it
via \texttt{run\_iaa(state, retriever=mock\_retriever)}, isolating chain logic from
Qdrant state. Integration tests (marked \texttt{@pytest.mark.integration}) bypass
the fixture and make real API calls.

\subsubsection{Corpus Ingestion Pipeline}

The ingestion pipeline (\texttt{ingest.py}) supports three source types: BPIC XES
process logs (parsed via pm4py), CUAD contract clauses (loaded from HuggingFace),
and tabular CSV/JSON event logs (via a column mapping config). All sources feed a
single Qdrant collection, with domain and process type stored as payload metadata.
Ingestion is idempotent: document IDs are deterministic MD5-based UUIDs
(\texttt{\_make\_id()}), so re-running ingestion on the same data skips already-present
documents. The \texttt{bootstrap\_corpus()} function provides a one-shot setup path;
the fastest bootstrap uses synthetic cases only and takes under 60 seconds including
embedding.

% ═════════════════════════════════════════════════════════════════════════════
\subsection{Challenges and Solutions}

\subsubsection{Gemini Quota Exhaustion}

\textbf{Problem.} The Gemini 2.0-Flash model on the free tier is limited to
20 requests per day. During demo dry-runs with multiple workflow cases (vendor
onboarding + claim review + fraud investigation), the quota was exhausted within
a single session, causing HTTP 429 errors from the Gemini API and IAA returning
HTTP 500 to the orchestrator.

\textbf{Solution.} Two changes were made. First, the \texttt{GEMINI\_MODEL} env var
was introduced so the model can be switched without code changes
(\texttt{gemini-2.0-flash-lite} offers 1,000 requests/day on the free tier at
reduced capability). Second, and more significantly, the fast-path bypass was
implemented: by returning a deterministic escalation decision for high/critical risk
cases without calling the LLM, approximately 40--60\% of demo cases no longer
consume quota. For a 3-case demo (one vendor, one claim, one fraud), the fraud case
always takes the fast path, leaving quota for the two standard-path cases.

\subsubsection{Qdrant File Lock in Local Mode}

\textbf{Problem.} During development, multiple modules (ingest script, retriever,
Streamlit app) each instantiated a \texttt{QdrantClient(path=./qdrant\_store)}.
The second instantiation in the same process raised the exception:
\begin{quote}\ttfamily Storage folder ./qdrant\_store is already accessed by another instance\end{quote}

\textbf{Solution.} \texttt{qdrant\_client\_singleton.py} implements a module-level
\texttt{\_client} global with lazy initialisation via \texttt{get\_qdrant\_client()}.
Every module that needs Qdrant access imports \texttt{get\_qdrant\_client} from
the singleton rather than constructing its own \texttt{QdrantClient}. The same
pattern is applied to the embedding model (\texttt{\_embedding\_model} global), which
also takes a significant hit on cold start (\textasciitilde{}90 seconds for the first
\texttt{SentenceTransformer("all-MiniLM-L6-v2")} call on a machine without the model cached).

\subsubsection{Contract Mismatch: context\_summary (Decision D1)}

\textbf{Problem.} IAA's \texttt{ProcessState} was originally designed with a
\texttt{context\_summary} field, expected to be populated by the Context Stitching
Agent with a plain-English summary of all retrieved documents. When the IAA chain
was first integrated into the AMP orchestrator, the orchestrator shim
(\texttt{csa\_to\_iaa()}) tried to pass \texttt{context\_summary} from CSA's output,
but CSA's \texttt{CsaOutput} model had no such field --- CSA returns structured JSON
(emails, documents, case\_history, calendar events), not a prose summary.
This caused a \texttt{KeyError} in the prompt assembly step.

\textbf{Solution.} The \texttt{context\_summary} field was demoted from required to
optional (\texttt{Optional[str] = None}) in \texttt{ProcessState}. The orchestrator
shim no longer attempts to pass it. The retriever's \texttt{build\_query()} function
includes \texttt{context\_summary} in the query only if it is non-null, so the query
degrades gracefully. This was documented as Decision D1 in the project decision log.

\subsubsection{BERT Model Cold Start}

\textbf{Problem.} The \texttt{sentence-transformers/all-MiniLM-L6-v2} model
(\textasciitilde{}80\,MB) is downloaded and loaded from disk on the first call to
\texttt{get\_embedding\_model()}. On a cold machine without the model in the
sentence-transformers cache, this takes 60--90 seconds and blocks the IAA server
from accepting requests.

\textbf{Solution.} The \texttt{/health} endpoint in \texttt{api.py} allows the
orchestrator to poll readiness before sending the first \texttt{/predict} request.
The \texttt{startup()} event handler in \texttt{api.py} pre-warms both the retriever
and LLM on server start, triggering the model download during the warm-up window.
The \texttt{start\_agents.sh} script uses a staggered startup sequence to give IAA
additional time before the orchestrator sends the first workflow.

\subsubsection{Fast-Path / Standard-Path Output Shape Consistency}

\textbf{Problem.} Early versions of the fast path returned a minimal
\texttt{ActionRecommendation} without \texttt{task\_list} and \texttt{workflow\_timeline}.
The downstream shim (\texttt{shims.py::iaa\_to\_amp()}) expected both fields to be
present and iterable, causing a \texttt{TypeError} when the orchestrator tried to
access \texttt{recommendation.task\_list} on a fast-path response.

\textbf{Solution.} Both fast-path constructors in \texttt{\_fast\_path()} now
explicitly pass \texttt{task\_list=[]}, \texttt{workflow\_timeline=[]} to the
\texttt{ActionRecommendation} constructor. The \texttt{ActionRecommendation} model
has \texttt{task\_list} and \texttt{workflow\_timeline} with \texttt{default\_factory=list},
but the fast path is explicit to avoid reliance on defaults. All override paths in
\texttt{apply\_overrides()} also thread through the original recommendation's
\texttt{task\_list} and \texttt{workflow\_timeline} using named constructor arguments.

\subsubsection{Three-Step Chain Latency}

\textbf{Problem.} The original chain design made three sequential LLM calls
(Contextualize $\to$ Anticipate $\to$ Recommend). On Gemini 2.5-Flash, each call
took approximately 8--10 seconds, giving a total latency of \textasciitilde{}26 seconds
per IAA invocation. This made the orchestrator's total workflow time exceed 30 seconds
on the standard path.

\textbf{Solution.} The three calls were merged into a single \texttt{MERGED\_SYSTEM} +
\texttt{MERGED\_USER} prompt that elicits all outputs (summary, predicted actions,
task list, workflow timeline, recommended action) in one JSON response. The JSON
parsing step was split into four separate parsers (for \texttt{predicted\_actions},
\texttt{task\_list}, \texttt{workflow\_timeline}, \texttt{recommended\_action}) applied
to the single LLM response. Total standard-path latency dropped from
\textasciitilde{}26\,s to \textasciitilde{}8--10\,s, and quota consumption dropped by 67\%.

% ═════════════════════════════════════════════════════════════════════════════
\subsection{Future Enhancements}

\subsubsection{Loop + Critic Pattern}

The current IAA makes a single LLM call and returns its output (modulo rule overrides).
A critic pass --- a second, shorter LLM call that receives the recommendation and
asks ``is this the right action given the case state?'' --- would reduce hallucinated
action labels and improve confidence calibration. The critic could also flag
cases where the recommended action contradicts the workflow timeline (e.g.\ recommending
a step that has already been completed). Implementation requires a second prompt template
and chain call in \texttt{chain.py}, returning a criticised version of the
\texttt{recommended\_action} with adjusted confidence. Complexity: medium. Risk:
doubles LLM latency on the standard path (to \textasciitilde{}16--20\,s); mitigated
by running the critic only for confidence below a configurable threshold.

\subsubsection{Retrieval Confidence Scoring}

Qdrant's \texttt{query\_points()} already returns cosine similarity scores for each
retrieved document (available as \texttt{r.score} in \texttt{retriever.py}). These
scores are currently discarded. Surfacing them in \texttt{PredictResponse} (e.g.\ as
a \texttt{retrieval\_score} field alongside \texttt{retrieved\_context\_ids}) would
allow the orchestrator and dashboard to weight IAA's output by retrieval quality:
a recommendation backed by highly similar historical cases (score $>$ 0.9) should be
trusted more than one backed by weakly similar cases (score $<$ 0.6). Complexity: low.
The \texttt{ActionRecommendation} schema would need one new field; no change to the
retrieval logic itself.

\subsubsection{Multi-Step Lookahead}

The current IAA predicts only the single next action. For well-known workflow types
(e.g.\ vendor onboarding: submit $\to$ verify $\to$ approve $\to$ onboard), the IAA
could predict the next $N$ steps, enabling DA to pre-assign work to the right agents
or humans for the entire known workflow path before the first step completes. This
would require a multi-turn LLM chain where each predicted action becomes the input
to the next prediction call. Complexity: high. Requires redesigning the chain to
accumulate a growing action sequence and stop when confidence falls below threshold.

\subsubsection{Automated Corpus Feedback Loop}

Currently the IAA corpus is static after initial bootstrapping. When the Delegation
Agent assigns a case and the case is eventually resolved, the outcome (what action
was actually taken and whether it led to a good resolution) is not fed back into the
corpus. Closing this loop would allow the RAG store to improve with every resolved
case: the \texttt{/ingest} endpoint already accepts runtime documents, so DA would
only need to POST the resolution event after case close. Complexity: medium.
Requires an event hook in DA (or a background job reading DA's audit log) and a
format converter from DA's \texttt{AssignmentDecision} to IAA's ingestion format.

\subsubsection{OpenAI Provider Completion}

\texttt{LLM\_PROVIDER=openai} is wired in \texttt{chain.py} and the env var
documentation, but has not been tested end-to-end through the full AMP workflow.
The \texttt{ChatOpenAI} factory branch exists (\texttt{chain.py:46}) and returns
\texttt{gpt-4o-mini} at \texttt{temperature=0.1}. What is missing is an integration
test (\texttt{@pytest.mark.integration}) that runs a full \texttt{run\_iaa()} call
with a real OpenAI key and verifies the response shape. Complexity: low. A single
integration test and a note in \texttt{.env.example} would complete the feature.

% ═════════════════════════════════════════════════════════════════════════════
\subsection{File Structure}

\begin{lstlisting}[style=ascii, caption={IAA directory tree with roles}]
IAA/
|
+-- agent/                        # Active production code
|   +-- __init__.py
|   +-- api.py                    # FastAPI REST server — port 8003 (active interface)
|   +-- chain.py                  # Fast path + merged LLM chain + rule overrides
|   +-- iaa.py                    # run_iaa() + iaa_node() LangGraph wrapper + audit log
|   +-- ingest.py                 # Corpus ingestion: BPIC / CUAD / CSV / synthetic
|   +-- main.py                   # ADK entry point (root_agent) — NOT used in AMP flow
|   +-- qdrant_client_singleton.py# Process-wide Qdrant + embedder singletons
|   +-- retriever.py              # Qdrant vector search + prompt formatting
|   +-- schemas.py                # All Pydantic v2 models (input, output, state)
|
+-- orchestrator/
|   +-- orchestrator.py           # Local 4-agent orchestrator — NOT the AMP LangGraph
|
+-- docs/                         # Reference documentation (read-only reference)
|   +-- data_ingestion.md         # Corpus strategy, embedding model, ingest commands
|   +-- deployment.md             # ADK web, REST API, Docker, env vars, known issues
|   +-- iaa_schema.json           # JSON Schema for IAA schemas (generated)
|   +-- iaa_schema.yaml           # YAML Schema for IAA schemas (generated)
|   +-- tech_stack.md             # LLM decisions, Qdrant singleton rationale, chain design
|   +-- testing.md                # Mock design, pytest commands, test structure
|
+-- data/                         # Runtime data (mostly gitignored)
|   +-- synthetic_cases.json      # 200 synthetic process cases for bootstrap corpus
|   +-- predictions.jsonl         # Audit log: one JSON line per IAA invocation
|
+-- qdrant_store/                 # Persisted Qdrant local store (gitignored at root)
|   +-- collection/
|   |   +-- iaa_corpus/           # Collection storage (SQLite + payload files)
|   +-- meta.json                 # Qdrant collection registry
|
+-- fixtures/                     # Test fixtures
|   +-- sample_case.json          # Single reference case for fixture-based tests
|
+-- test_iaa.py                   # 77 unit tests across 8 test classes
+-- conftest.py                   # autouse LLM mock — zero API calls in tests
+-- iaa_demo.py                   # CLI demo: list cases, single case, high-risk, workflow
+-- requirements.txt              # Python dependencies
+-- CLAUDE.md                     # Agent collaboration rules (checked in)
\end{lstlisting}

\textbf{Active code} (called in the production AMP flow):
\texttt{api.py}, \texttt{chain.py}, \texttt{iaa.py}, \texttt{retriever.py},
\texttt{schemas.py}, \texttt{qdrant\_client\_singleton.py}, \texttt{ingest.py}.

\textbf{Reference / standalone only} (not called by the AMP orchestrator):
\texttt{agent/main.py} (ADK interface), \texttt{orchestrator/orchestrator.py}
(IAA's own local orchestrator, used by \texttt{/workflow} endpoint and CLI demo only).

\textbf{Test infrastructure}:
\texttt{test\_iaa.py}, \texttt{conftest.py}, \texttt{fixtures/sample\_case.json}.

\end{document}
